\documentclass[9pt,a4paper]{extarticle}
\usepackage{parskip}

% ======================
% Packages
% ======================
\usepackage[utf8]{inputenc}
\usepackage[T1]{fontenc}
\usepackage[brazil]{babel}
\usepackage{amsmath, amssymb}
\usepackage{amsthm}
\usepackage{geometry}
\usepackage{listings}
\usepackage{xcolor}
\usepackage{hyperref}
\usepackage{booktabs}
\usepackage{array}
\usepackage{tabularx}
\usepackage{enumitem}
\usepackage{microtype}
\usepackage{csquotes}
\usepackage{natbib}

\definecolor{important}{RGB}{255, 100, 100}
\definecolor{notice}{RGB}{0, 102, 204}

\newcommand{\impt}[1]{\textcolor{important}{#1}}
\newcommand{\ntc}[1]{\textcolor{notice}{#1}}

\setlist{nosep}
\geometry{margin=2cm}

% ======================
% Listings setup (Python)
% ======================
\definecolor{codebg}{rgb}{0.95,0.95,0.95}
\lstset{
    language=Python,
    backgroundcolor=\color{codebg},
    basicstyle=\ttfamily\footnotesize,
    keywordstyle=\color{blue},
    commentstyle=\color{gray},
    stringstyle=\color{red!70!black},
    showstringspaces=false,
    frame=single,
    breaklines=true,
    inputencoding=utf8
}

% ======================
% Document
% ======================
\begin{document}

\title{Introdução \footnote{Material baseado em \citep{chalmers1993ciencia}}}
\author{BCC502 -- Metodologia Científica para Ciência da Computação}
\date{\today}
\maketitle

\tableofcontents

% ==============================================================
\section{Indutivismo}

\subsection{Revisão da posição indutivista}

\begin{itemize}
    \item \ntc{Tese central}: o conhecimento científico é \emph{conhecimento provado}, derivado dos fatos da experiência por meio de \impt{generalização indutiva}.
    \item A ciência parte de proposições de observação (afirmações singulares) e ascende a leis e teorias (afirmações universais) pela indução.
    \item A partir de leis e teorias, previsões e explicações são obtidas por \ntc{dedução}.
    \item A indução é o que distingue a ciência de meras opiniões: a generalização só é legítima quando há (i) muitas observações, (ii) sob ampla variedade de condições, (iii) sem contraexemplo.
\end{itemize}

\paragraph{Esquema do método indutivo}

\[
\underbrace{\text{Fatos adquiridos por observação}}_{\text{afirmações singulares}}
\;\xrightarrow{\text{indução}}\;
\underbrace{\text{Leis e teorias}}_{\text{afirmações universais}}
\;\xrightarrow{\text{dedução}}\;
\underbrace{\text{Previsões e explicações}}_{\text{afirmações singulares}}
\]

\subsection{Problemas do indutivismo}

\paragraph{O problema da indução (Hume)}
\begin{itemize}
    \item Não há justificação \emph{lógica} para a inferência indutiva: a verdade das premissas não garante a verdade da conclusão.
    \item Justificar a indução pelo seu sucesso passado é circular --- usa-se indução para justificar a indução.
    \item Portanto, \impt{leis e teorias científicas não são, em sentido estrito, ``provadas'' pelos fatos}.
\end{itemize}

\paragraph{O problema da base observacional}
\begin{itemize}
    \item \impt{A observação não é neutra}: o que se observa depende do conhecimento prévio, das expectativas e do quadro teórico do observador.
    \item Dois observadores diante do mesmo objeto não necessariamente registram a mesma proposição de observação.
    \item Logo, a ideia de uma base de fatos puros, anterior e independente da teoria, é insustentável.
\end{itemize}

\paragraph{Imprecisão das condições indutivas}
\begin{itemize}
    \item Quantas observações são ``muitas''? Quais condições contam como ``variadas''?
    \item Sem uma teoria que selecione as variáveis relevantes, a exigência de variedade é vazia.
\end{itemize}

\texttt{Conclusão: a ciência não pode ser caracterizada como conhecimento provado a partir de fatos. É preciso outra concepção de método científico.}

% ==============================================================
\section{Falsificacionismo}
% ==============================================================

\subsection{A virada popperiana}

\begin{itemize}
    \item \ntc{Tese central (Popper)}: o que distingue a ciência não é a possibilidade de \emph{provar} suas teorias, mas a de \impt{refutá-las}.
    \item A assimetria lógica entre verificação e refutação é a chave: nenhum número finito de observações prova uma afirmação universal, mas \impt{uma única observação contrária basta para refutá-la}.
    \item \ntc{Critério de demarcação}: uma teoria é científica se, e somente se, for \impt{falsificável} --- isto é, se existirem proposições de observação possíveis que, se verdadeiras, a refutariam.
\end{itemize}

\paragraph{Forma lógica da falsificação}

\begin{itemize}
    \item Premissa 1 (teoria): \(\forall x\,(A(x) \rightarrow B(x))\)
    \item Premissa 2 (observação): \(A(a) \land \lnot B(a)\)
    \item Conclusão: a teoria é falsa.
\end{itemize}

Esta inferência é \ntc{dedutivamente válida} (\textit{modus tollens}), ao contrário da inferência indutiva.

\subsection{Graus de falsificabilidade}

\begin{itemize}
    \item Teorias melhores são as \impt{mais falsificáveis} --- aquelas que proíbem mais.
    \item Uma teoria é tanto mais informativa quanto mais arriscadas forem suas previsões.
    \item Exemplo: ``Os planetas se movem em órbitas elípticas em torno do Sol'' é mais falsificável (e mais científica) do que ``Os planetas se movem em curvas fechadas''.
\end{itemize}

\paragraph{Ciência como conjecturas e refutações}

\begin{itemize}
    \item As teorias são \impt{conjecturas ousadas}, não generalizações cautelosas.
    \item A ciência progride por tentativas de refutação: propõe-se uma hipótese, deduzem-se previsões arriscadas, testam-se.
    \item \ntc{Corroboração} não é prova: uma teoria que resiste a testes severos é apenas \emph{provisoriamente aceita}.
\end{itemize}

\subsection{Limites do falsificacionismo}

\paragraph{A tese de Duhem--Quine}
\begin{itemize}
    \item Teorias não são testadas isoladamente: junto com a hipótese principal estão sempre hipóteses auxiliares, condições iniciais e teorias instrumentais.
    \item Diante de uma observação contrária, é sempre possível \impt{salvar a teoria revisando uma hipótese auxiliar}.
    \item Portanto, a falsificação conclusiva é, na prática, impossível.
\end{itemize}

\paragraph{O problema histórico}
\begin{itemize}
    \item Cientistas raramente abandonam uma teoria diante do primeiro contraexemplo.
    \item Muitas teorias hoje bem-sucedidas (Newton, Bohr) \emph{nasceram refutadas} por evidências disponíveis.
    \item Um falsificacionismo ingênuo descartaria essas teorias --- e estaria errado em fazê-lo.
\end{itemize}

\texttt{Estes limites motivam concepções mais estruturadas, em que a unidade de avaliação não é a teoria isolada, mas conjuntos articulados de teorias ao longo do tempo.}

% ==============================================================
\section{Programas de Pesquisa Científica (Lakatos)}
% ==============================================================

\subsection{A proposta de Lakatos}

\begin{itemize}
    \item Lakatos busca preservar o que há de correto no falsificacionismo sem ignorar a história real da ciência.
    \item \ntc{Tese central}: a unidade de avaliação não é uma teoria isolada, mas um \impt{programa de pesquisa} --- uma sequência historicamente desenvolvida de teorias.
    \item Um programa não é refutado por uma anomalia: é \emph{abandonado} quando deixa de produzir resultados, em favor de um programa rival mais promissor.
\end{itemize}

\subsection{Estrutura de um programa de pesquisa}

\paragraph{Núcleo duro (\textit{hard core})}
\begin{itemize}
    \item Conjunto de hipóteses fundamentais que definem o programa.
    \item Tomado por decisão metodológica como \impt{irrefutável}: anomalias nunca são direcionadas ao núcleo.
    \item Exemplo: as três leis do movimento e a lei da gravitação no programa newtoniano.
\end{itemize}

\paragraph{Cinturão protetor (\textit{protective belt})}
\begin{itemize}
    \item Conjunto de hipóteses auxiliares, condições iniciais e teorias observacionais que cercam o núcleo.
    \item \impt{Recebe o impacto das refutações}: é onde se fazem ajustes diante das anomalias.
    \item É modificado, refinado e até substituído ao longo do desenvolvimento do programa.
\end{itemize}

\paragraph{Heurísticas}
\begin{itemize}
    \item \ntc{Heurística negativa}: regra metodológica que proíbe dirigir o \textit{modus tollens} contra o núcleo.
    \item \ntc{Heurística positiva}: programa de trabalho que orienta como desenvolver e modificar o cinturão protetor --- quais modelos construir, quais aproximações relaxar, em que ordem atacar os problemas.
\end{itemize}

\subsection{Programas progressivos e degenerativos}

\begin{itemize}
    \item \impt{Progressivo}: modificações no cinturão protetor são \emph{teoricamente progressivas} (preveem fatos novos) e \emph{empiricamente progressivas} (parte dessas previsões é corroborada).
    \item \impt{Degenerativo}: modificações são \textit{ad hoc} --- servem apenas para acomodar anomalias já conhecidas, sem prever nada novo.
    \item A escolha entre programas rivais não se faz por um experimento crucial isolado, mas pela \ntc{avaliação comparativa, ao longo do tempo}, de sua produtividade.
\end{itemize}

\texttt{Um programa degenerativo pode, em princípio, voltar a ser progressivo. Por isso, abandoná-lo é uma decisão metodológica, não uma refutação lógica.}

% ==============================================================
\section{Paradigmas de Pesquisa (Kuhn)}
% ==============================================================

\subsection{A proposta de Kuhn}

\begin{itemize}
    \item Kuhn parte de uma análise \impt{histórica e sociológica} da prática científica, não de uma reconstrução lógica.
    \item \ntc{Tese central}: a ciência madura é organizada em torno de \impt{paradigmas} --- realizações científicas universalmente reconhecidas que, por algum tempo, fornecem problemas e soluções modelares para uma comunidade.
    \item O desenvolvimento científico não é cumulativo: alterna períodos de estabilidade com rupturas profundas.
\end{itemize}

\subsection{O ciclo kuhniano}

\paragraph{Pré-ciência}
\begin{itemize}
    \item Atividade desorganizada, com escolas rivais e desacordo sobre fundamentos.
    \item Não há ainda um paradigma compartilhado.
\end{itemize}

\paragraph{Ciência normal}
\begin{itemize}
    \item Pesquisa firmemente baseada num paradigma aceito pela comunidade.
    \item Consiste em \impt{resolução de quebra-cabeças} (\textit{puzzle-solving}): aplicar e articular o paradigma, não testá-lo.
    \item Anomalias não são vistas como refutações, mas como problemas a serem resolvidos dentro do paradigma.
\end{itemize}

\paragraph{Crise}
\begin{itemize}
    \item Acúmulo de anomalias persistentes e resistentes a soluções dentro do paradigma.
    \item Confiança da comunidade enfraquece; tentativas \textit{ad hoc} se multiplicam; surgem candidatos a paradigma alternativo.
\end{itemize}

\paragraph{Revolução científica}
\begin{itemize}
    \item Substituição do paradigma antigo por um novo.
    \item Não é uma escolha puramente lógica: envolve fatores estéticos, pragmáticos, sociológicos e geracionais.
\end{itemize}

\paragraph{Nova ciência normal} 

Reinicia-se o ciclo sob o novo paradigma.

\subsection{O conceito de paradigma}

\begin{itemize}
    \item Um paradigma é mais do que uma teoria: inclui leis, princípios, técnicas experimentais, padrões de aplicação, valores epistêmicos e exemplares (problemas-modelo).
    \item Define, para a comunidade, \impt{o que conta como problema legítimo, método aceitável e solução adequada}.
    \item Funciona simultaneamente como \ntc{matriz disciplinar} (compromissos compartilhados) e \ntc{exemplar} (soluções concretas que servem de modelo).
\end{itemize}

\subsection{Incomensurabilidade}

\begin{itemize}
    \item \impt{Tese da incomensurabilidade}: paradigmas rivais não compartilham inteiramente vocabulário, problemas ou critérios de avaliação.
    \item Não há uma linguagem observacional neutra que permita compará-los ponto a ponto.
    \item Adotar um novo paradigma se assemelha mais a uma \ntc{conversão} do que a uma decisão por evidência decisiva.
    \item Isto não implica irracionalismo: a comunidade científica decide com base em critérios partilhados (precisão, alcance, simplicidade, fecundidade), mas estes critérios não determinam univocamente a escolha.
\end{itemize}

\texttt{Lakatos e Kuhn convergem em rejeitar o falsificacionismo ingênuo, mas divergem na ênfase: Lakatos preserva uma racionalidade metodológica avaliando programas pela sua progressividade; Kuhn enfatiza o papel da comunidade e dos fatores extra-lógicos nas mudanças de paradigma.}

% ==============================================================
\section{Racionalismo e Relativismo}
% ==============================================================

A discussão sobre Lakatos e Kuhn nos leva a uma questão mais profunda: \texttt{existe um critério universal e atemporal que permita decidir, racionalmente, entre teorias científicas rivais?} As respostas a essa pergunta organizam-se em dois polos --- o racionalismo e o relativismo.

\subsection{Racionalismo}

\begin{itemize}
    \item \ntc{Tese central}: existe um \impt{critério universal, atemporal e impessoal} que permite avaliar os méritos relativos de teorias rivais.
    \item Esse critério é \emph{independente do contexto} --- não depende da época, da cultura, dos interesses ou dos valores da comunidade que avalia.
    \item Filósofos e cientistas, enquanto racionais, devem aceitar ou rejeitar teorias conforme esse critério se aplica.
\end{itemize}

\paragraph{Implicações do racionalismo}
\begin{itemize}
    \item \impt{Demarcação clara}: o critério separa o que é ciência do que não é.
    \item \impt{Progresso objetivo}: faz sentido dizer que a ciência avança em direção à verdade (ou a teorias melhores) num sentido absoluto.
    \item As decisões científicas, quando bem feitas, não dependem de fatores sociológicos, psicológicos ou históricos.
\end{itemize}

\paragraph{Exemplos de critérios racionalistas}
\begin{itemize}
    \item \ntc{Indutivismo}: aceite teorias que tenham sido suficientemente confirmadas pela indução a partir dos fatos.
    \item \ntc{Falsificacionismo (Popper)}: aceite (provisoriamente) teorias que sejam altamente falsificáveis e que tenham resistido a tentativas severas de refutação.
\end{itemize}

\subsection{Relativismo}

\begin{itemize}
    \item \ntc{Tese central}: \impt{não existe um critério universal e atemporal} para avaliar teorias. O que conta como conhecimento, método válido e boa teoria \emph{depende} do indivíduo ou da comunidade que avalia.
    \item Padrões de avaliação são \emph{internos} a uma cultura, época, paradigma ou tradição --- e variam com eles.
    \item Não há um ponto de vista privilegiado, fora de toda tradição, a partir do qual julgar as tradições.
\end{itemize}

\paragraph{Implicações do relativismo}
\begin{itemize}
    \item \impt{Não há demarcação absoluta} entre ciência e não-ciência; a distinção é convencional e historicamente situada.
    \item Falar em ``progresso'' só faz sentido \emph{dentro} de uma tradição --- não há progresso absoluto rumo à verdade.
    \item Mudanças teóricas profundas envolvem necessariamente fatores que vão além da lógica e da evidência: valores comunitários, formação, persuasão, contexto histórico.
\end{itemize}

\paragraph{Uma distinção importante}

Relativismo não equivale a dizer que ``vale tudo''. Dentro de uma tradição, há critérios rigorosos, debate racional e decisões fundamentadas. O que o relativista nega é que esses critérios sejam \impt{universais} --- isto é, que valham igualmente para toda comunidade, em toda época.

\subsection{Resumo do contraste}

\begin{center}
\begin{tabularx}{\linewidth}{@{} l X X @{}}
\toprule
 & \textbf{Racionalismo} & \textbf{Relativismo} \\
\midrule
Critério de avaliação & Universal, atemporal, impessoal & Interno a uma comunidade ou tradição \\
Demarcação ciência/não-ciência & Possível e bem definida & Convencional, dependente do contexto \\
Progresso científico & Objetivo, rumo à verdade & Definido apenas dentro de uma tradição \\
Mudança teórica & Decidida por lógica e evidência & Envolve também fatores sociológicos e históricos \\
\bottomrule
\end{tabularx}
\end{center}

\texttt{Onde se situam Lakatos e Kuhn nesse espectro?}

% ==============================================================
\section{Lakatos e Kuhn diante do Racionalismo e do Relativismo}
% ==============================================================

\subsection{Lakatos: um racionalismo sofisticado}

\begin{itemize}
    \item Lakatos é \impt{explicitamente racionalista}: sua metodologia dos programas de pesquisa pretende oferecer um critério universal para avaliar a ciência.
    \item Esse critério, no entanto, \ntc{não se aplica a teorias isoladas}, mas a programas de pesquisa avaliados ao longo do tempo.
    \item O critério é a distinção entre programas \impt{progressivos} e \impt{degenerativos}: programas que preveem fatos novos e os corroboram são racionalmente preferíveis àqueles que apenas absorvem anomalias \textit{ad hoc}.
\end{itemize}

\paragraph{Por que é racionalismo}
\begin{itemize}
    \item O critério (progressivo vs.\ degenerativo) é proposto como \emph{universal}: aplica-se em qualquer época, qualquer área, qualquer comunidade.
    \item A escolha entre programas rivais é, em princípio, uma decisão \emph{racional}, fundamentada na produtividade comparada.
    \item Fatores sociológicos podem explicar por que de fato cientistas se comportam de certa maneira, mas não justificam \emph{racionalmente} essa escolha.
\end{itemize}

\paragraph{Por que é \emph{sofisticado}}
\begin{itemize}
    \item Diferentemente do falsificacionismo ingênuo, Lakatos reconhece que (i) teorias nascem cercadas de anomalias, (ii) a falsificação conclusiva é impossível, (iii) a avaliação requer tempo e perspectiva histórica.
    \item Ele admite uma \ntc{tolerância racional} ao erro: um programa pode ser legitimamente mantido mesmo com anomalias, desde que continue a produzir previsões bem-sucedidas.
\end{itemize}

\subsection{Kuhn: leituras racionalistas e leituras relativistas}

A posição de Kuhn é mais ambígua e foi objeto de longa controvérsia. Há elementos de sua obra que puxam para o relativismo e outros que recuam dele.

\paragraph{Elementos relativistas em Kuhn}
\begin{itemize}
    \item A \impt{tese da incomensurabilidade}: paradigmas rivais não compartilham inteiramente vocabulário, problemas ou critérios. Logo, não há um ponto de vista neutro a partir do qual compará-los.
    \item Adotar um novo paradigma se assemelha a uma \ntc{conversão}: envolve fatores estéticos, geracionais e sociológicos, não apenas evidência.
    \item O que conta como problema legítimo, método aceitável e boa solução é definido \emph{internamente} pelo paradigma. Não há um critério externo de cientificidade.
    \item Não faz sentido falar em ``aproximação da verdade'' ao longo do desenvolvimento histórico da ciência.
\end{itemize}

\paragraph{Elementos racionalistas em Kuhn}
\begin{itemize}
    \item Kuhn insiste que a escolha entre paradigmas \emph{não é arbitrária}.
    \item Há \impt{valores epistêmicos compartilhados} pela comunidade científica: precisão, alcance, simplicidade, consistência, fecundidade.
    \item Esses valores não determinam univocamente a escolha (cientistas razoáveis podem, com base neles, preferir paradigmas diferentes), mas a restringem.
    \item Kuhn rejeitou expressamente o rótulo de relativista em escritos posteriores e procurou esclarecer que a incomensurabilidade é \emph{local} e \emph{linguística}, não uma incomunicabilidade total.
\end{itemize}

\paragraph{A controvérsia}
\begin{itemize}
    \item Críticos (notadamente o próprio Lakatos) leram Kuhn como reduzindo a mudança científica a \emph{psicologia das multidões} --- um relativismo disfarçado.
    \item Defensores leem Kuhn como descrevendo a racionalidade científica \emph{realmente praticada}, que é mais rica e situada do que a racionalidade abstrata dos filósofos.
\end{itemize}

\subsection{Resumo do contraste}

\begin{center}
\begin{tabularx}{\linewidth}{@{} l X X @{}}
\toprule
 & \textbf{Lakatos} & \textbf{Kuhn} \\
\midrule
Posição declarada & Racionalismo (sofisticado) & Rejeita o rótulo, mas leitura ambígua \\
Unidade de avaliação & Programa de pesquisa & Paradigma / comunidade científica \\
Critério de avaliação & Progressivo vs.\ degenerativo (universal) & Valores epistêmicos compartilhados (locais) \\
Fatores extra-lógicos & Explicam, mas não justificam & Constitutivos da prática científica \\
Progresso & Objetivo: programas progressivos avançam & Há progresso \emph{a partir} de onde se está, não \emph{em direção a} uma verdade \\
\bottomrule
\end{tabularx}
\end{center}

\texttt{A questão de fundo permanece em aberto: é possível uma teoria da ciência que seja, ao mesmo tempo, fiel à história real da prática científica e capaz de justificar racionalmente as escolhas teóricas?}

% ==============================================================
\section{Objetivismo (O que é uma teoria?)} 
% ==============================================================

\subsection{Posição central}

\begin{itemize}
    \item \ntc{Tese central}: o conhecimento científico tem propriedades que \impt{não dependem das crenças, intenções ou estados mentais dos indivíduos} que o produzem ou utilizam.
    \item Uma vez formulado, o conhecimento ``ganha vida própria'': pode ter consequências, implicações e propriedades que ninguém ainda percebeu.
    \item O foco filosófico recai sobre o conhecimento como \emph{produto} --- proposições, teorias, problemas, métodos --- e não sobre os \emph{processos mentais} que levaram a ele.
\end{itemize}

\paragraph{Contraste fundamental}

\begin{itemize}
    \item \ntc{Individualismo (ou subjetivismo)}: o conhecimento é antes de tudo um \impt{estado mental} de indivíduos. Falar em conhecimento científico é falar no que cientistas \emph{acreditam}, \emph{percebem} ou \emph{entendem}.
    \item \ntc{Objetivismo}: o conhecimento é uma \impt{estrutura}, com propriedades próprias, situada fora das mentes individuais. Cientistas trabalham \emph{sobre} essa estrutura, mas ela não se reduz a eles.
\end{itemize}

\texttt{Exemplo: a teoria de Maxwell tem como consequência a existência de ondas eletromagnéticas. Essa consequência estava ``na teoria'' antes que Maxwell --- ou qualquer outra pessoa --- a percebesse. Hertz a descobriu décadas depois.}

\subsection{Argumentos a favor do objetivismo}

\paragraph{Consequências não percebidas}
\begin{itemize}
    \item Teorias têm implicações lógicas que excedem em muito o que seus autores tinham em mente.
    \item Essas implicações \impt{existem na estrutura da teoria}, à espera de serem extraídas.
    \item Se o conhecimento fosse apenas estado mental, falar de ``consequências ainda não percebidas'' seria contraditório.
\end{itemize}

\paragraph{Conhecimento como bem público}
\begin{itemize}
    \item Teorias podem ser registradas em livros, criticadas, testadas, refinadas e ensinadas \emph{independentemente} do estado psicológico de quem as formulou.
    \item Uma teoria continua existindo, com suas propriedades, mesmo após a morte de seu autor.
    \item A ciência é, nesse sentido, uma \impt{atividade pública} cujos resultados transcendem indivíduos.
\end{itemize}

\paragraph{Problemas situados em teorias, não em mentes}
\begin{itemize}
    \item Quando uma teoria é proposta, certos \ntc{problemas surgem objetivamente} dela --- tensões internas, anomalias, lacunas --- sem que ninguém precise tê-los notado.
    \item Esses problemas são propriedades da estrutura da teoria, não invenções psicológicas.
    \item A pesquisa científica é, em larga medida, o trabalho de localizar e atacar esses problemas.
\end{itemize}

\subsection{Implicações para a filosofia da ciência}

\paragraph{O foco da análise muda}
\begin{itemize}
    \item Para o individualista, perguntas centrais são: \emph{em que o cientista acreditava?}, \emph{como ele chegou a essa ideia?}, \emph{que processo psicológico produziu a descoberta?}
    \item Para o objetivista, as perguntas centrais são: \emph{quais são as propriedades lógicas e estruturais da teoria?}, \emph{que problemas ela enfrenta?}, \emph{que consequências ela tem?}, \emph{como se compara a teorias rivais?}
    \item \impt{A psicologia da descoberta é separada da avaliação lógica do produto.}
\end{itemize}

\paragraph{Relação com as outras concepções estudadas}
\begin{itemize}
    \item \ntc{Falsificacionismo de Popper}: fortemente objetivista. As teorias são avaliadas por sua estrutura lógica (grau de falsificabilidade, conteúdo empírico), não pela confiança subjetiva de quem as defende.
    \item \ntc{Programas de pesquisa de Lakatos}: também objetivista. Um programa é progressivo ou degenerativo por propriedades \emph{suas} --- da sequência de teorias e suas previsões ---, não pelo entusiasmo da comunidade que o adota.
    \item \ntc{Paradigmas de Kuhn}: aqui surge uma tensão. Como o paradigma envolve crenças compartilhadas, formação, valores da comunidade, a posição de Kuhn é frequentemente lida como menos objetivista. A escolha entre paradigmas envolve elementos que parecem inseparáveis dos estados mentais da comunidade científica.
\end{itemize}

\subsection{Objeções ao objetivismo}

\paragraph{Conhecimento sem sujeitos?}
\begin{itemize}
    \item Críticos argumentam que faz pouco sentido falar em conhecimento sem alguém que \emph{conheça}.
    \item Teorias só se realizam como conhecimento quando alguém as compreende, aplica, contesta.
    \item Resposta objetivista: o conhecimento \emph{em potência} existe na estrutura, ainda que precise de mentes para ser ativado --- assim como uma partitura é música em potência mesmo quando ninguém a executa.
\end{itemize}

\paragraph{A história sugere o oposto}
\begin{itemize}
    \item Pesquisas em história e sociologia da ciência mostram que decisões científicas \emph{de fato} envolvem fortes componentes psicológicos, institucionais e sociais.
    \item Um objetivismo radical correria o risco de descrever uma ciência idealizada, distante da praticada.
    \item Resposta objetivista: descrever como cientistas \emph{de fato} agem (descritivo) é diferente de avaliar a \emph{qualidade epistêmica} do que produzem (normativo). O objetivismo é uma tese sobre o segundo, não sobre o primeiro.
\end{itemize}

\texttt{O objetivismo oferece uma resposta possível à tensão entre racionalismo e relativismo: separar a avaliação lógica das teorias --- que é objetiva --- da história sociológica de sua aceitação --- que envolve fatores contingentes. Saber se essa separação é sustentável é uma questão filosófica em aberto.}

\bibliographystyle{apalike}   
\bibliography{references}

\end{document}

